\exercisesection{同调维数}

\begin{enumerate}
\def\labelenumi{\arabic{enumi})}
\tightlist
\item
  设 \(u: \mathbf{A} \to \mathbf{B}\) 是一个环同态，M 是一个 B-模。证明不等式
\end{enumerate}

\[
\mathrm{dp} _ {\mathbf {A}} (\mathbf {M}) \leqslant \mathrm{dp} _ {\mathbf {B}} (\mathbf {M}) + \mathrm{dp} _ {\mathbf {A}} (\mathbf {B} _ {s}).
\]

（通过对 dpB(M) 归纳推理，或利用第 188 页习题 6 的谱序列。）

\begin{enumerate}
\def\labelenumi{\arabic{enumi})}
\setcounter{enumi}{1}
\item
  假设环 A 是诺特的；设 M 是一个有限生成的 A-模，其投射维数为 n。证明 k-模 Ext \(_{A}^{n}\) (M, A) 是非零的。
\item
  设 \(n\) 为一个整数 \(\geqslant 1\) 。证明 \(\mathrm{dh}(\mathbf{M}_n(\mathbf{A})) = \mathrm{dh}(\mathbf{A})\) （参见 VIII，§ 4）。
\item
  设 M 是一个 A-模。我们称 M 的内射维数（记为 \(\mathrm{di}_{\mathbf{A}}(\mathbf{M})\) ）为 M 的内射分解的长度在 \(\overline{\mathbf{Z}}\) 中的下确界。
\end{enumerate}

如果 n 是一个整数 \(\geqslant\) 0，证明以下条件是等价的：

\(\alpha)\mathrm{di}_{\mathbf{A}}(\mathbf{M})\leqslant n;\)

β) \(\operatorname{Ext}_{\mathbf{A}}^{r}(\mathbf{N},\mathbf{M}) = 0\) 对每个 A-模 N 及每个整数 \(r > n\) ；

\(\gamma)\ \mathrm{Ext}_{\mathbf{A}}^{n+1}(\mathbf{N},\mathbf{M}) = 0\) 对于每个 A-模 \(\mathbf{N}\) ；

δ) \(\operatorname{Ext}_{\mathbf{A}}^{n+1}(\mathbf{N},\mathbf{M})=0\) 对于每个循环 A-模 N；

ε) 对于每个正合列 0 → M → I⁰ → \(\cdots\) → Iⁿ⁻¹ → N → 0，其中诸 Iᵏ 都是内射的，N 是内射的。 5) 设 (Mᵢ)ᵢ∈F 是一族 A-模。

\begin{enumerate}
\def\labelenumi{\alph{enumi})}
\item
  证明 \(\mathrm{di}_{\mathbf{A}}(\prod \mathbf{M}_i) = \sup \mathrm{di}_{\mathbf{A}}(\mathbf{M}_i)\) 。
\item
  若 A 是诺特的，证明 \(\mathrm{di}_{\mathbf{A}}(\bigoplus \mathbf{M}_i) = \sup \mathrm{di}_{\mathbf{A}}(\mathbf{M}_i)\) 。
\item
  若 \(\mathrm{di}_{\mathbf{A}}(\bigoplus_{i\in E}\mathbf{M}_{i}) = \sup_{i\notin E}\mathrm{di}_{\mathbf{A}}(\mathbf{M}_{i})\) 对每个 A-模族 \((\mathbf{M}_i)_{i\in E}\) 成立，则 A 是诺特的（利用第 170 页习题 21）。
\end{enumerate}

\begin{enumerate}
\def\labelenumi{\arabic{enumi})}
\setcounter{enumi}{5}
\item
  设 I 为 A 的左理想组成的集合。证明 \(\mathrm{dh}(\mathbf{A}) = \sup \mathrm{dp}_{\mathbf{A}}(\mathbf{A}/\mathfrak{a})\) 。
\item
  设 M 为一个 A-模。我们称 M 的平坦维数（记为 \(dpl_{A}\) (M)）为 M 的平坦分解的长度在 \(\overline{Z}\) 中的下确界。
\end{enumerate}

\begin{enumerate}
\def\labelenumi{\alph{enumi})}
\tightlist
\item
  如果 n 是一个整数 \(\geqslant 0\) ，证明以下条件是等价的：
\end{enumerate}

\(\alpha)\mathrm{dpl}_{\mathbf{A}}(\mathbf{M})\leqslant n.\)

\(\beta)\ \mathrm{Tor}_{r}^{\mathbf{A}}(\mathbf{N},\mathbf{M})=0\ \text{对每个右 A-模 N 及每个整数 }r>n.\)

\(\gamma)\ \mathrm{Tor}_{n+1}^{\mathbf{A}}(\mathbf{N},\mathbf{M})=0\) 对每个右 A-模 N。

\(\delta)\ \mathrm{Tor}_{n+1}^{\mathbf{A}}(\mathbf{N},\mathbf{M})=0\) 对每个单生成右 A-模 N。

ε) 对每个正合列 \(0 \rightarrow K \rightarrow P_{n-1} \rightarrow \cdots \rightarrow P_{0} \rightarrow M \rightarrow 0\) ，其中诸 \(P_{k}\) 是平坦的，则 K 是平坦的。

\begin{enumerate}
\def\labelenumi{\alph{enumi})}
\setcounter{enumi}{1}
\item
  证明有 \(\mathrm{dpl}_{\mathbf{A}}(\mathbf{M}) \leqslant \mathrm{dp}_{\mathbf{A}}(\mathbf{M})\) ；若 A 是诺特环且 M 是有限生成的，则等号成立。
\item
  设 \((\mathbf{N}_{\alpha})_{\alpha \in I}\) 为 A-模的一个归纳系统，N 为其归纳极限。证明不等式
\end{enumerate}

\[
\mathrm{dpl} _ {A} (N) \leqslant \sup _ {\alpha \in I} \mathrm{dpl} _ {A} (N _ {\alpha}).
\]

\begin{enumerate}
\def\labelenumi{\arabic{enumi})}
\setcounter{enumi}{7}
\tightlist
\item
  我们称 A 的 tor 维数（记为 td(A)）为在 \(\overline{\mathbf{Z}}\) 中、由满足下述条件的整数 \(n\) 组成的集合的上确界：存在一个左 A-模 N 和一个右 A-模 M，使得 \(\operatorname{Tor}_{n}^{\mathbf{A}}(\mathbf{M},\mathbf{N})\neq0\) 。
\end{enumerate}

\begin{enumerate}
\def\labelenumi{\alph{enumi})}
\tightlist
\item
  设 \(n\) 为一个整数 \(\geqslant 0\) 。证明以下条件是等价的：
\end{enumerate}

\(\alpha)\ \mathrm{td}(\mathbf{A}) \leqslant n.\)

β) 对于每个 A-模 M，我们有 \(\mathrm{dpl}_{\mathbf{A}}(\mathbf{M}) \leqslant n\) （习题 7）。

\(\gamma)\) 对于每个单生成 A-模 M，我们有 \(\mathrm{dpl}_{\mathbf{A}}(\mathbf{M}) \leqslant n\) 。

\begin{enumerate}
\def\labelenumi{\alph{enumi})}
\setcounter{enumi}{1}
\tightlist
\item
  证明我们有 td(A) = td(A°) 且 td(A) ≤ dh(A)；若 A 是左诺特的，则等号成立。
\end{enumerate}

¶ 9) a) 证明以下条件是等价的：

α) A 的每个有限生成（左）理想都是投射 A-模。

β) 投射 A-模的每个有限生成子模都是投射的。

\(\gamma)\) 我们有 td(A) \(\leqslant\) 1（习题8），且环A是凝聚的（X，第181页，习题11）。

δ) 对任意集合 I，A¹ 的每个子模都是平坦的。

\begin{enumerate}
\def\labelenumi{\alph{enumi})}
\setcounter{enumi}{1}
\tightlist
\item
  设环 A 满足 a) 中的等价条件。证明：每个投射 A-模都同构于 A 的有限生成理想的直和。（先处理有限生成模的情形；然后对具有可数生成元族的模用归纳法处理，注意到投射 A-模 P 的每个元素都包含在 P 的一个有限生成直和项中。最后从 Kaplansky 定理（II，第 183 页，习题 2）推出一般情形。）
\end{enumerate}

\begin{enumerate}
\def\labelenumi{\arabic{enumi})}
\setcounter{enumi}{9}
\tightlist
\item
  假设 \(\mathbf{A}\) 是整环；设 \(\mathbf{K}\) 为其分式域。我们说一个理想 \(\mathfrak{a}\) （\(\mathbf{A}\) 的）是 \(invertible\) （可逆的），如果存在元素 \(x_{1},\ldots ,x_{n}\) （\(\mathbf{K}\) 中的），使得 \(x_{i}\mathfrak{a}\subset \mathbf{A}\) （\(1\leqslant i\leqslant n\) ）且 \(\sum x_i\mathfrak{a} = \mathbf{A}\) 。
\end{enumerate}

\begin{enumerate}
\def\labelenumi{\alph{enumi})}
\item
  证明每个可逆理想都是有限生成的。
\item
  证明 A 的非零理想是投射 A-模当且仅当它是可逆的。
\item
  设 a 是 A 的一个可逆理想，D 是一个可除 A-模。证明 \(\operatorname{Ext}_A^i(A/a, D) = 0\) （\(i \geqslant 1\) ）。
\end{enumerate}

\begin{enumerate}
\def\labelenumi{\arabic{enumi})}
\setcounter{enumi}{10}
\tightlist
\item
  \begin{enumerate}
  \def\labelenumii{\alph{enumii})}
  \tightlist
  \item
    假设环 A 是一个整环。证明习题 9 的条件 \(\alpha\) 到 \(\delta\) 仍然等价于以下条件：
  \end{enumerate}
\end{enumerate}

ε) 每个有限生成的挠自由 A-模都是投射的。

ζ) 每个无挠 A-模都是平坦的。

η) 平坦 A-模的每个子模都是平坦的。

（利用第 169 页习题 13。）

满足这些条件的整环称为普吕弗环。

\begin{enumerate}
\def\labelenumi{\alph{enumi})}
\setcounter{enumi}{1}
\tightlist
\item
  设 B 是一个整环，它是一个递增滤过的普吕弗子环族的并。证明 B 是普吕弗的。
\end{enumerate}

* c) 证明代数整数环是普吕弗环。 *

\begin{enumerate}
\def\labelenumi{\arabic{enumi})}
\setcounter{enumi}{11}
\tightlist
\item
  假设环A是整环。证明以下条件是等价的：
\end{enumerate}

\(\alpha)\ \mathrm{dh}(\mathbf{A}) \leqslant 1.\)

β) 每个可除A-模都是内射的。

γ) 环 A 是普吕弗环（习题 11）且是诺特环。

（为证明 α) 蕴含 β)，使用习题 10。）

于是我们称 A 为戴德金环。

\begin{enumerate}
\def\labelenumi{\arabic{enumi})}
\setcounter{enumi}{12}
\tightlist
\item
  假设A是左且右诺特的。证明以下条件是等价的：
\end{enumerate}

\(\alpha)\ \mathrm{dh}(\mathbf{A}) \leqslant 2.\)

β) 对于所有整数 \(p, q\) 以及每个 A-线性映射 \(u: \mathbf{A}^p \to \mathbf{A}^q\) ， \(u\) 的核是投射 A-模。

γ) 每个有限生成A-模的对偶都是投射的。

\begin{enumerate}
\def\labelenumi{\arabic{enumi})}
\setcounter{enumi}{13}
\tightlist
\item
  设 x 是 A 中的一个左可消的元素。
\end{enumerate}

\begin{enumerate}
\def\labelenumi{\alph{enumi})}
\item
  证明该理想 \(x\mathbf{A}\) 是双侧的，当且仅当存在 \(\sigma\) ——环 \(\mathbf{A}\) 的一个自同态——使得 \(ax = x\sigma(a)\) 对所有 \(a \in \mathbf{A}\) 。
\item
  我们现在假设 a) 的条件成立。设 M 是一个 A-模，使得 xM = 0。证明等式 dp\_A(M) = dp\_{A/xA}(M) + 1（可以利用第 188 页习题 6，d) 的谱序列）。
\item
  我们假设 \(x\) 属于 \(\mathbf{A}\) 的根，并且 \(\mathbf{A}\) 是诺特的。设 \(\mathbf{N}\) 是一个有限生成的 \(\mathbf{A}\) -模，使得比为 \(x\) 的位似变换在 \(\mathbf{N}\) 上是单射的。证明我们有 \(\mathrm{dp}_{\mathbf{A}}(\mathbf{N}) = \mathrm{dp}_{\mathbf{A}/\mathbf{A}x}(\mathbf{N}/x\mathbf{N})\) 。d) 在 \(c\) ) 的条件下，证明有 \(\mathrm{dh}(\mathbf{A}) = \mathrm{dh}(\mathbf{A}/\mathbf{A}x) + 1\) 。
\end{enumerate}

\begin{enumerate}
\def\labelenumi{\arabic{enumi})}
\setcounter{enumi}{14}
\tightlist
\item
  设 σ 是环 A 的一个自同构。
\end{enumerate}

\begin{enumerate}
\def\labelenumi{\alph{enumi})}
\tightlist
\item
  我们用 \(\mathbf{A}_{\sigma}[\mathbf{X}]\) 表示在 VIII，§ 1，第 3 号中定义的环。证明
\end{enumerate}

\[
\mathrm{dh} (\mathbf {A} _ {\sigma} [ \mathbf {X} ]) = \mathrm{dh} (\mathbf {A}) + 1.
\]

（不等式 dh(A \(_{\sigma}\) {[}X{]}) \(\geqslant\) dh(A) + 1 由习题 14， \(b\) ) 得出；对于相反的不等式，可借鉴 X，第 142 页，引理 3， \(a\) 的证明。）

\begin{enumerate}
\def\labelenumi{\alph{enumi})}
\setcounter{enumi}{1}
\tightlist
\item
  我们用 \(\mathbf{A}_{\sigma}[[\mathbf{X}]]\) 表示在 IV，§ 4，习题 8 中定义的环。证明：若 A 是诺特环，则我们有 \(\mathrm{dh}(\mathbf{A}_{\sigma}[[\mathbf{X}]))=\mathrm{dh}(\mathbf{A})+1\) （利用习题 14）。
\end{enumerate}

\begin{enumerate}
\def\labelenumi{\arabic{enumi})}
\setcounter{enumi}{15}
\tightlist
\item
  设 A 是一个右诺特且左诺特的局部环，m 为其极大理想；我们用 \(k_{s}\) （相应地， \(k_{d}\) ）表示左（相应地，右）A-模 A/m。设 M 是一个 Tor \(_{n+1}^{A}\) 有限生成左模，并设 n 为一个整数。a) 证明 dp \(_{A}\) (M) \(\leqslant\) n 当且仅当 Tor \(_{n+1}^{A}\) ( \(k_{d}\) , M) = 0（利用第 185 页习题 4）。
\end{enumerate}

\begin{enumerate}
\def\labelenumi{\alph{enumi})}
\setcounter{enumi}{1}
\tightlist
\item
  证明有 \(\mathrm{dh(A)} = \mathrm{dp}_{\mathbf{A}}(k_s)\) ；要使 \(\mathrm{dh(A)} \leqslant n\) 成立，其充分必要条件是有
\end{enumerate}

\[
\operatorname{Tor} _ {n + 1} ^ {\mathbf {A}} \left(k _ {d}, k _ {s}\right) = 0.
\]

\begin{enumerate}
\def\labelenumi{\alph{enumi})}
\setcounter{enumi}{2}
\tightlist
\item
  设 \(x\) 是 \(\mathbf{A}\) 的中心的一个元素，属于 \(m\) ，使得比值为 \(x\) 的位似映射在 \(M\) 中是单射的。证明我们有 \(\mathrm{dp}_{\mathrm{A}}(\mathrm{M}/x\mathrm{M}) = \mathrm{dp}_{\mathrm{A}}(\mathrm{M}) + 1\) （另参见习题 14）。
\end{enumerate}

\begin{enumerate}
\def\labelenumi{\arabic{enumi})}
\setcounter{enumi}{16}
\tightlist
\item
  设 M 是一个 A-模，I 是一个良序集， \((\mathbf{M}_{\alpha})_{\alpha \in I}\) 是 M 的子模的一个递增族，使得 \(M = \bigcup M_{\alpha}\) 。我们设 \(M_{\alpha}^{\prime} = \bigcup M_{\beta}\) （\(\alpha \in I\) ）。
\end{enumerate}

\begin{enumerate}
\def\labelenumi{\alph{enumi})}
\item
  证明不等式 dpA(M) ≤ sup dpA(Mα/Mα')。
\item
  假设存在一个整数 \(n\) 使得 \(\mathrm{dp}_{\mathbf{A}}(\mathbf{M}_{\alpha}') \leqslant n\) 对所有 \(\alpha \in I\) 成立。证明我们有
\end{enumerate}

\[
\mathrm{dp} _ {\mathbf {A}} (\mathbf {M}) \leqslant n + 1.
\]

\begin{enumerate}
\def\labelenumi{\arabic{enumi})}
\setcounter{enumi}{17}
\tightlist
\item
  \begin{enumerate}
  \def\labelenumii{\alph{enumii})}
  \tightlist
  \item
    设 \((\mathbf{M}_{\alpha})_{\alpha \in I}\) 是相对于可数集 I 的 A-模的一个归纳系统，M 为其极限， \(n\) 为一个整数 \(\geqslant 0\) 。证明如果 \(\mathrm{dp}_{\mathbf{A}}(\mathbf{M}_{\alpha}) \leqslant n\) 对所有 \(\alpha \in I\) 成立，则 \(\mathrm{dp}_{\mathbf{A}}(\mathbf{M}) \leqslant n + 1\) （化归到 \(n = 0\) 且 \(I = N\) 的情形，然后写出正合列 \(0 \to \bigoplus_{n} \mathbf{M}_{n} \to \bigoplus_{n} \mathbf{M}_{n} \to \mathbf{M} \to 0\) ）。
  \end{enumerate}
\end{enumerate}

\begin{enumerate}
\def\labelenumi{\alph{enumi})}
\setcounter{enumi}{1}
\item
  设 P 是一个平坦 A-模，它具有一个表示 \(\mathbf{A}^{(\mathbb{S})}\to\mathbf{A}^{(\mathbb{T})}\to\mathbf{P}\to0\) ，其中集合 S 是可数的。由 a) 推出 dp \(_{\mathbb{A}}(\mathbb{P})\leqslant1\) 。
\item
  假设环 A 的每个左理想都由一个可数元素集合生成。证明：对于由可数族元素生成的每个 A-模 M，有 dpA(M) ≤ dplA(M) + 1。由此推出不等式 dh(A) ≤ td(A) + 1 以及 \textbar dh(A) - dh(A°)\textbar{} ≤ 1。（利用第 181 页习题 12、第 202 页习题 8，以及 b）。）
\end{enumerate}

¶ 19) a) 设 n 为整数 ≥ 0，设 (Mα)α ∈ I 是 A-模的一个归纳系，并设 M 为其归纳极限。假设 Card (I) ≤ Nn（E，III，第 87 页，习题 10），且存在整数 r 使得对每个 α ∈ I 有 dpA(Mα) ≤ r。证明 dpA(M) ≤ r + n + 1（通过对 n 归纳论证，利用习题 18，a) 和习题 17）。

\begin{enumerate}
\def\labelenumi{\alph{enumi})}
\setcounter{enumi}{1}
\tightlist
\item
  我们假设 A 的每个左理想都由一个基数为 \(\leqslant \aleph_{n}\) 的元素族生成。证明我们有 dp \(_{A}\) (M) \(\leqslant\) dpl \(_{A}\) (M) + n + 1，对每个由基数为 \(\leqslant \aleph_{n}\) 的元素族生成的 A-模 M 成立。由此推出不等式
\end{enumerate}

\[
\mathrm{dh} (\mathbf {A}) \leqslant \mathrm{td} (\mathbf {A}) + n + 1 \quad \text { and } \quad | \mathrm{dh} (\mathbf {A}) - \mathrm{dh} (\dot {\mathbf {A}} ^ {\circ}) | \leqslant n + 1.
\]

¶ 20) 设 R 为主理想环，K 为其分式域。我们用 B 表示 M₂(K) 中由形如 \(\begin{pmatrix} a & 0 \\ x & y \end{pmatrix}\) （\(a \in \mathbb{R}, x, y \in \mathbb{K}\) ）的矩阵组成的子环。

\begin{enumerate}
\def\labelenumi{\alph{enumi})}
\item
  证明B是左诺特的，但不是右诺特的。
\item
  证明 \(\mathrm{dh(B)} = 1\) 且 \(\mathrm{dh(B^{\circ})} = 2\) 。
\end{enumerate}

\begin{enumerate}
\def\labelenumi{\arabic{enumi})}
\setcounter{enumi}{20}
\tightlist
\item
  设 r 为 A 的根。证明以下条件等价：
\end{enumerate}

α) A/r 是半单的，并且对 r 中元素的每一个序列 \((a_{n})_{n\geq1}\) ，存在一个整数 n 使得 \(a_{1}\ldots a_{n}=0\) 。

β) 每个A-模都有投射覆盖。

\(\gamma)\) 对于每个 A-模 M，我们有 \(\mathrm{dpl}_{\mathbf{A}}(\mathbf{M}) = \mathrm{dp}_{\mathbf{A}}(\mathbf{M})(\mathbf{X}, p. 202, \text{exercise } 7)\) 。

δ) 对 A-模的每一个归纳系统 \((\mathbf{M}_{\alpha}, f_{\beta\alpha})_{\alpha \in I}\) ，我们有：

\[
\mathrm{dp} _ {\mathrm{A}} (\varinjlim \mathrm{M} _ {\alpha}) \leqslant \sup _ {\alpha \in I} \mathrm{dp} _ {\mathrm{A}} (\mathrm{M} _ {\alpha}).
\]

ε) 每个平坦的A-模都是投射的。

φ) A 的单生成右理想的每个递减序列都是平稳的。

(对于 \(\alpha\) ）与 \(\beta\) ）的等价性，参见 VIII，§ 8，习题 26。为证明 \(\beta\) ）蕴含 \(\gamma\) ），注意 M 具有一个极小投射分解 P，且 \(P_n/rP_n\) 同构于 Tor \(_n^A\) (A/r, M)。为证明 \(\varepsilon) \Rightarrow \varphi\) ），设 \((l_n)_{n \geq 1}\) 是单生成右理想的一个递减序列，使得

\[
\mathrm{I} _ {n} = a _ {1} \dots a _ {n} \mathrm{A}, \quad \text { with } \quad a _ {i} \in \mathrm{A}.
\]

在自由模 \(\mathbf{A}^{(\mathbf{N})}\) 中，考虑子模 \(\mathbf{N}_p\) ，它由 \(e_1 - a_1 e_2, \ldots, e_p - a_p e_{p+1}\) 生成。从 \(\varepsilon\) ）推出 \(\mathbf{N} = \bigcup \mathbf{N}_p\) 是 \(\mathbf{A}^{(\mathbf{N})}\) 中的直和项；由此得出结论： \(\mathrm{l}_{n+1} = \mathrm{l}_n\) （当 \(n\) 足够大时）。

对于蕴含关系 \(\varphi) \Rightarrow \alpha)\) ，使用 VIII，§ 9，习题 12。）

\begin{enumerate}
\def\labelenumi{\arabic{enumi})}
\setcounter{enumi}{21}
\tightlist
\item
  证明以下条件是等价的：
\end{enumerate}

(α) 每个具有有限投射维数的有限表示 A-模都是投射的。

(β) 每个单射 A-同态 \(A^{p} \rightarrow A^{q}\) 都有一个收缩。

(γ) 对每一个不同于 A 的有限生成右理想 a，存在 A 中一个非零元素 x，使得 xa = 0。

(δ) 每个有限表示的单右 A-模都同构于 A 的一个（极小）右理想。

（注意到条件 \(\gamma\) ) 等价于条件 \(\beta\) ），当 \(p = 1\) 时。）

¶ 23) 证明以下条件等价：

(α) 每个具有有限投射维数的 A-模都是投射的。

(β) 环 A 满足习题 21 的等价条件，并且对于每个不同于 A 的有限生成右理想 a，存在 A 的一个元素 \(x \neq 0\) 使得 \(xa = 0\) 。

(γ) 环 A 满足习题 21 中的等价条件，且每个单 A-模都是某个内射 A-模的商。

（证明 \(\alpha\) ) 蕴含 A 的单生成右理想的每个递减序列都是平稳的，这可通过改编习题 21 的证明得到。然后由习题 22 推出 \(\alpha\) ）与 \(\beta\) ）的等价性。为看出 \(\alpha\) ）蕴含 \(\gamma\) ），考虑单模 M 的一个投射覆盖 P、一个包含 P 的内射模 I，以及 I 的一个投射覆盖 Q；定义一个从 Q 到 P 的 A-同态，并由此推出一个从 I 到 M 的满射。为证明 \(\gamma\) ）蕴含 \(\alpha\) )，证明一个满足 dp \(_{A}\) (M) \(\leqslant\) 1 的 A-模 M 是投射的，这通过考虑一个投射覆盖 \(p: P \to M\) 以及一个从 Ker \(p\) 到单模上的满射来完成。）

\begin{enumerate}
\def\labelenumi{\arabic{enumi})}
\setcounter{enumi}{23}
\tightlist
\item
  我们用 dhf(A)（相应地，dif(A)、tdf(A)）表示使该维数有限的 A-模的投射（相应地，内射、平坦）维数的上确界。若 A 具有有限同调维数，则 dhf(A) = dif(A) = dh(A) 且 tdf(A) = td(A)。若 td(A) \textless{} ∞，我们有 tdf(A) = td(A)。 a) 证明 tdf(A) ≤ dif(A°)；若 A 是左诺特的，则等号成立（利用 X，第 189 页，习题 3）。
\end{enumerate}

\begin{enumerate}
\def\labelenumi{\alph{enumi})}
\setcounter{enumi}{1}
\item
  假设 A 是左诺特的。证明我们有 dhf(A) ≤ di\_A(A\_s)。若 di\_A(A\_s) 与 dif(A) 都是有限的，证明它们相等（利用第 202 页习题 2）。
\item
  设 G 是一个不等于单位元的有限群，并设 A 为环 \(\mathbf{Z}^{(G)}\) 。证明我们有
\end{enumerate}

\[
\mathrm{dh} (\mathbf {A}) = + \infty , \mathrm{dhf} (\mathbf {A}) = \mathrm{dif} (\mathbf {A}) = \mathrm{tdf} (\mathbf {A}) = 1, \mathrm{di} _ {\mathbf {A}} (\mathbf {A} _ {s}) = + \infty
\]

（参见第 193 页习题 20）。

\begin{enumerate}
\def\labelenumi{\alph{enumi})}
\setcounter{enumi}{3}
\tightlist
\item
  假设环 A 是自内射的（参见 X，第 171 页，习题 26），且不是半单的。证明每个非投射 A-模都有无限投射维数和无限内射维数。由此推出 dh(A) = +∞，dhf(A) = dif(A) = tdf(A) = di\_A(A\_s) = 0。
\end{enumerate}

\begin{enumerate}
\def\labelenumi{\arabic{enumi})}
\setcounter{enumi}{24}
\tightlist
\item
  我们用 \(\mathbf{A}^e\) 表示 \(k\) -代数 \(\mathbf{A} \otimes_k \mathbf{A}^*\) ，并且我们把 \(\mathbf{A}\) 视为一个左 \(\mathbf{A}^e\) -模。我们设
\end{enumerate}

\[
\mathrm{da} _ {k} (\mathrm{A}) = \mathrm{dp} _ {\mathrm{A} ^ {e}} (\mathrm{A}).
\]

如果 n 是一个整数 \(\geqslant0\) ，则不等式 \(\mathrm{da}_{k}(\mathrm{A})\leqslant n\) 等价于 \(\mathrm{H}^{n+1}(\mathrm{A},\mathrm{M})\) 的消没（对每个 (A, A)-双模 M）（X，第 195 页，习题 26）。我们有 \(\mathrm{da}_{k}(\mathrm{A})=\mathrm{da}_{k}(\mathrm{A}^{\circ})\) 。

\begin{enumerate}
\def\labelenumi{\alph{enumi})}
\item
  证明有 \(\mathrm{da}_k(\mathbf{M}_n(k)) = 0\) 且 \(\mathrm{da}_k(k[\mathbf{X}_1, \ldots, \mathbf{X}_n]) = n\) （对所有 \(n \geqslant 0\) ）。
\item
  若 \(k\) 是一个域，则 \(\mathrm{da}_{k}(\mathbf{A}) = 0\) 当且仅当 \(k\) -代数 \(\mathbf{A}\) 是绝对半单的（参见 VIII，§ 11，n° 3，定理 2）。
\item
  设 M 是一个自由 \(k\) -模。证明我们有 \(\mathrm{da}_{k}(\mathbf{T}_{k}(\mathbf{M})) = 1\) （若 \((e_{i})_{i \in I}\) 是 M 的一组基，证明元素 \((e_{i} \otimes 1 - 1 \otimes e_{i})\) 构成 \(\mathbf{A}^{e}\) 的左理想的一组基，该左理想是乘法 \(\mathbf{A} \otimes_{k} \mathbf{A}^{-} \to \mathbf{A}\) 的核，其中 \(\mathbf{A} = \mathbf{T}_{k}(\mathbf{M})\) ）。
\item
  现在我们假设 \(k\) -模 A 是投射的。设 \(\rho: k \to k'\) 是交换环的一个同态；证明有 \(\mathrm{da}_k(\mathbf{A} \otimes_k k') \leqslant \mathrm{da}_k(\mathbf{A})\) 。若 \(\rho\) 是单射且 \(\rho(k)\) 是 \(k\) -模 \(k'\) 的一个直和项，则有等式（利用 X，第 195 页，习题 26，c)）。
\item
  设 B 是第二个 k-代数（结合的且有单位元），在 k 上投射。证明：
\end{enumerate}

\[
\mathrm{da} _ {k} (\mathbf {A} \otimes_ {k} \mathbf {B}) \leqslant \mathrm{da} _ {k} (\mathbf {A}) + \mathrm{da} _ {k} (\mathbf {B}).
\]

此外，若 k 是一个域，且 A 和 B 在 k 上是有限维的，则等号成立（利用 X，第 195 页，习题 26，d)）。

\begin{enumerate}
\def\labelenumi{\alph{enumi})}
\setcounter{enumi}{5}
\tightlist
\item
  假设 k 是一个域。证明不等式 \(\mathrm{dh}(\mathbf{A}) \leqslant \mathrm{da}_{k}(\mathbf{A})\) 与 \(\mathrm{dh}(\mathbf{A}^{\circ}) \leqslant \mathrm{da}_{k}(\mathbf{A})\) （参见 X，第 196 页，习题 26，e)）。
\end{enumerate}

\begin{enumerate}
\def\labelenumi{\arabic{enumi})}
\setcounter{enumi}{25}
\tightlist
\item
  设 A 为一个增广 \(k\) -代数（X，第 196 页，习题 27）；假设存在一个同态 \(\rho: \mathrm{A} \to \mathrm{A}^e\) ，它满足 X，第 196 页，习题 27，c) 的条件。证明 \(\mathrm{da}_k(\mathbf{A}) = \mathrm{dp}_\mathbf{A}(k)\) ；此外若 \(k\) 是一个域，则 \(\mathrm{da}_k(\mathbf{A}) = \mathrm{dp}_\mathbf{A}(k) = \mathrm{dh}(\mathbf{A}) = \mathrm{dh}(\mathbf{A}^\circ)\) 。（利用第 196 页习题 27 的 b) 和 c)，以及第 196 页习题 26 的 e)。）特别地，若 V 是一个非零 \(k\) -向量空间，我们有 \(\mathrm{dh}\,\mathbf{T}_k(\mathbf{V}) = 1\) 。
\end{enumerate}

* 27) 我们假设 \(k\) 是一个域。设 \(g\) 是域 \(k\) 上的一个李代数，维数为 \(n\) ，U 为其包络代数。证明 \(dhU = n\) 。（利用第 196 页习题 26、习题 27，d)，以及第 197 页习题 28，c)。）*

¶ 28) 设 G 是一个群。称 G 的上同调维数，记作 dc(G)，为 \(\mathbf{Z}^{(G)}\) -模 Z 的投射维数。若 \(n\) 是一个整数，则我们有 dc(G) \(\leqslant n\) 当且仅当 H \(^{q}\) (G, M)=0 （对每个 \(\mathbf{Z}^{(G)}\) -模 M 和每个 \(q > n\) ）。

\begin{enumerate}
\def\labelenumi{\alph{enumi})}
\item
  证明有 \(\mathrm{da}_{\mathbf{Z}}(\mathbf{Z}^{(G)}) = \mathrm{dc}(G)\) （利用第 196 页习题 26 和习题 27，d)）。
\item
  若 G 是一个不约化为单位元的自由群，则 dc(G) = 1。
\item
  设 H 是 G 的子群；证明不等式 dc(H) ≤ dc(G)。若 dc(G) \textless{} ∞ 且 H 在 G 中具有有限指数，证明 dc(H) = dc(G)（证明当 q = dc(G) 时，第 191 页习题 14 中的同态
\end{enumerate}

\[
t ^ {q}: \mathrm{H} ^ {q} (\mathrm{H}, \mathrm{M}) \rightarrow \mathrm{H} ^ {q} (\mathrm{G}, \mathrm{M})
\]

是满射）。

\begin{enumerate}
\def\labelenumi{\alph{enumi})}
\setcounter{enumi}{3}
\item
  证明一个不约化为单位元的有限群的上同调维数是无穷。由此推出，若 dc(G) \textless{} ∞，则 G 中除单位元外的每个元素都有无穷阶（此时称 G 是无挠的）。
\item
  若 G 是无挠的且 H 是 G 中有限指数的子群，证明 dc(H) = dc(G)。（若 (P, p) 是 \(\mathbf{Z}^{(H)}\) -模 Z 的一个投射分解，且若 (G : H) = n，则在分次 \(\mathbf{Z}^{(H)}\) -模 \(\mathbf{T}_{\mathbf{Z}}^{n}(\mathbf{P})\) 上定义一个 \(\mathbf{Z}^{(G)}\) -复形结构，使得 \((\mathbf{T}_{\mathbf{Z}}^{n}(\mathbf{P}), p^{\otimes n})\) 是 \(\mathbf{Z}^{(G)}\) -模 Z 的一个投射分解。）
\item
  若 H 是 G 的正规子群，证明我们有
\end{enumerate}

\[
\mathrm{dc} (\mathbf {G}) \leqslant \mathrm{dc} (\mathbf {H}) + \mathrm{dc} (\mathbf {G} / \mathbf {H}).
\]

